\documentclass[12pt]{article}
\usepackage{lmodern}
\usepackage{amssymb,amsmath}
\usepackage{ifxetex,ifluatex}
\ifnum 0\ifxetex 1\fi\ifluatex 1\fi=0 % if pdftex
  \usepackage[T1]{fontenc}
  \usepackage[utf8]{inputenc}
  \usepackage{textcomp} % provides euro and other symbols
\else % if luatex or xelatex
  \usepackage{unicode-math}
  \defaultfontfeatures{Ligatures=TeX,Scale=MatchLowercase}
\fi
% use upquote if available, for straight quotes in verbatim environments
\IfFileExists{upquote.sty}{\usepackage{upquote}}{}
% use microtype if available
\IfFileExists{microtype.sty}{%
\usepackage[]{microtype}
\UseMicrotypeSet[protrusion]{basicmath} % disable protrusion for tt fonts
}{}
\IfFileExists{parskip.sty}{%
\usepackage{parskip}
}{% else
\setlength{\parindent}{0pt}
\setlength{\parskip}{6pt plus 2pt minus 1pt}
}
\usepackage{hyperref}
\hypersetup{
            pdfborder={0 0 0},
            breaklinks=true}
\setlength{\emergencystretch}{3em}  % prevent overfull lines
\providecommand{\tightlist}{%
  \setlength{\itemsep}{0pt}\setlength{\parskip}{0pt}}
\setcounter{secnumdepth}{0}





\usepackage[english]{babel}
\usepackage[utf8]{inputenc}
\usepackage{fancyhdr}




\newcommand{\myparagraph}[1]{\paragraph{#1}\mbox{}\\}
\usepackage{geometry}
\geometry{margin=1.5in}
\renewcommand{\baselinestretch}{1.12} 
\usepackage{graphicx}

 


\pagestyle{fancy}
\fancyhf{}
\rhead{{\small Bertini Baldassini Folco et Louzieni Yannis}}
\lhead{\small \leftmark}
\rfoot{Page \thepage}
\usepackage{multicol}
\usepackage{times}

\setlength{\parindent}{2em}


%{\usepackage[compact]{titlesec}
%\titlespacing{\section}{0pt}{2ex}{1ex}
%\titlespacing{\subsection}{0pt}{1ex}{0ex}
%\titlespacing{\subsubsection}{0pt}{0.5ex}{0ex}}

\date{}

\begin{document}
\begin{titlepage}
	\begin{center}
		\vspace*{1cm}
	
		{\Large \textbf{Rapport Orientation et Insertion Professionelle}}
		
		\vspace{0.5cm}
		Mèmes, Mémétique
		
		\vspace{1.5cm}
		
		\textbf{Bertini Baldassini Folco et Louzieni Yannis}
		

		\vfill


 		\includegraphics[width=7cm]{serveimage.png}
		\vspace{1.2cm}
		
		Un document rédigé pour le cours de\\
		Diana Ruiz Pino
		
		\vspace{1.2cm}
		
		Sorbonne Université\\
		Orientation et Insertion Professionnelle\\
		Paris, France\\
		Date
		
	\end{center}
\newpage
\tableofcontents

\end{titlepage}

\section{Introduction}
L'intéresse commun pour la mémétique et les mots de Richard Dawkins ont formé le binome auteur de ce document. Il se propose de analyser et comprendre l'interaction des mèmes sur internet et en particulier sur les réseaux sociaux. Pour arriver à cet objectif nous avons développé nos dossiers personnels en se rapprochant à l'argument, puis nous avons parcouru le sujet de la mémétique selon la méthode de l'OLMDP pour en comprendre l'ensemble et conclure sur la spécificité des réseaux.

\section{Mèmes}

\subsection{Objet}

\begin{quote}
	\emph{Cet objet a été choisi car loin de l'image que l'on s'en fait d'un
		sujet trivial et uniquement à visée humoristique, il s'agit de l'élément
		porteur d'une théorie qui apporte une réponse pertinente et originale
		Aux fonctionnements des comportements culturels. En effet, la mémétique
		étend la théorie de l'évolution a des phénomènes non plus naturels mais
		culturels. En proposant l'existence des mèmes qui évoluent de manière
		biologique on arrive à un modèle cohérent sans avoir besoins d'ajouter
		pléthore de nouvelles notions. D'autant plus que la mémétique se base
		sur des modèles très solides évolution darwinienne, entre autres. 	\cite{claidiere:tel-00431055} Aussi ayant remarqué la récurrence de même étapes dans l'évolution de		plusieurs mèmes (internet), j'ai été attiré par la capacité à expliquer
		à l'aide de règle simple les schémas de variation et de sélection de ces
		mèmes.}
\end{quote}

\emph{\textbf{Mots clés:}} Même -- mémétique -- néodarwinisme - gène
égoïste - théorie évolutionniste étendue

\subsubsection{Constituants}

Le mème est l'unité de hérédité culturelle. Il s'agit d'un concept
abstrait, impossible à manipuler, il est possible que d'en analyser les
principes de fonctionnement.\cite{definition} Il est constitué par le réplicateur, une
information codifiée qui produit copies exactes de soi même avec des
occasionnelles copies incorrectes ou ``mutations''. Les réplicateurs
mieux adaptés à faire des copies de soi même deviennent plus nombreux,
contrairement aux moins dotés.

Dans le monde des mèmes, les éventuels défauts de réplication
(mutations) provoquent une variété du même mème dans le pool mémétique
et les différents variantes pourront être en competition entre elles. La
compétition se déroule sur le plan de la phénotypique, c'est à dire sur
le comportement général d'un mème dans la société. Le destin d'un mème
dépend de sa transmission parmi les individus et les groupes ethniques. \cite{gene}
Avec le temps, les mèmes augmentent ou diminuent de fréquence dans le
pool mémétique selon des caractères phénotypiques que les représentent.
De plus, les mèmes coopérants ont une probabilité majeure d'émerger.

Les mèmes ont une capacité et fidélité élevées de transmission des
données, même si dans certains cas ils peuvent rester inactifs pendant
longtemps.

Exemple: le jeu du passe-mot avec une technique plus avancée, puis avec
un dessin (2 types d'information, analogique ou numérique/digitale) \cite{dio}

Un ``mémeplexe'' est un complexe de mèmes qui, même s'ils ne sont pas
nécessairement capables de survivre tous seuls, sont capables de
survivre en présence des autres mèmes du mémeplexe.


\subsubsection{Caractéristiques}

Il peut s'agir d'élément triviaux (blague récurrente, symbole
populaire), ou d'un objet d'une importance capitale (une langue, une
éthique). Le mème a une durée qui peut être très variable de quelques
mois à plusieurs millénaires.

Selon cette hypothèse toute culture peut être décomposée en ces sous
unités transmissibles entre les membres. En fait tous les éléments
constitutifs de notre culture sont des mèmes. Une chanson, par exemple,
est un mème. Si je chante une chanson, je serais entendu par quelqu'un,
et la chanson que j'ai chantée sera dupliquée dans un nouveau cerveau,
et ainsi de suite. L'imitation permet la réplication.

La réplication est donc l'essence du même, une idée qui ne se réplique
plus, n'a plus cours, et donc «~meure». Pour qu'un mème subsiste il faut
donc qu'il induise chez «~l'hôte~» un comportement susceptible d'être
imité, pour que celui-ci le propage à d'autres.

Selon la théorie de Dawkins, si le gène est le support de l'information
génétique, le même est celui de la culture. Ainsi donc, les mèmes
évoluent de la selon des lois analogue à celles qui régissent
l'évolution des gènes. \cite{dion:hal-00765151}

Ces lois sont~:

\begin{itemize}
\tightlist
\item
  Le principe de variation~: les instances sont differentes au sein de
  la «~population~» (c'est difference son dûes au mode de réplication.
  Pour filer la métaphore, il peut s'agir d'une mauvaise imitation de
  l'air ou des paroles de la chanson, ou bien d'une parodie, d'un
  détournement)
\item
  la selection du plus apte~: les instances qui auront une forme (ou un
  contenu) favorisant la réplication, seront sélectionnés les autres
  seront oubliés. En d'autres termes les mèmes les plus imitables
  perdurent.
\end{itemize}

Cependant plus un même est répliqué plus il accumule des variations~: il
peut donc donner naissance à plusieurs versions, voire des mèmes
dérivées. Ils peuvent aussi se combiner entre eux.

On notera aussi que plusieurs même en apparence différents peuvent
véhiculer une idée identique.

On peut se demander comment les mèmes peuvent-ils se répendre~? quel est
leur support~? En vérité il n'est pas aisé de dégager un support
privilégié, comme il serait difficile d'en dégager un pour la culture.
Néanmoins il est certain que les mèmes se véhiculent par tous les moyens
de communications, notamment inernet, d'où le terme «~Internet memes~».\cite{chronique}


\subsubsection{Méthodes d'obtentions}

Les méthodes d'obtention d'un mème sont nombreuses et souvent ne
dépendent pas d'un individu ou d'une volonté spécifique.

Il faut analyser les deux phases: la création et la diffusion. La
création d'un mème est la conséquence de la mutation d'un mème déjà
existant. Ce mème peut se manifester dans la phénotypique avec la
variation de la façon de prononcer une lettre ou un mot, un tic, une
expression, la répétition d'une mot, l'acquis d'une nouvelle habitude. \cite{origine}
Le changement d'un individu ou d'une société par un seul mème est
minimal et invisible, il faut que le pool mémétique change d'un 0,02\%
afin que la différence soit visible (l'utilisation fréquente d'un mot
par un individu). Des phénomènes plus grands, comme un mouvement
artistique, sont provoqués par une multitude de mèmes qui collaborent
entre eux et qui sont partagés par plusieurs individus. Pour cela il
faut en analyser la diffusion. Les mèmes conceptuels, ceux qui
deviennent partie du pool des connaissances globales, sont souvent crées
par les acteurs d'un certain secteur. Par exemple, les mèmes contenants
les théories scientifiques sont le produit du travail des scientifiques.

Les mèmes se diffusent par le partage d'idées, par imitation, par copie,
par éducation, par influence, par diffusion sur les réseaux sociaux,
etc. La survie d'un mème est déterminée par sa diffusion entre les gens
et les contenants d'informations. Un livre est souvent la représentation
d'un système de stockage de mèmes. Il détient le façons de parler, de se
comporter, de s'habiller et de penser d'une époque spécifique. Le mème
peut être totalement inactif dans la société mais rien n'empêche qu'il
puisse à nouveau se diffuser grâce à quelqu'un qui lit le livre et en
assume des mèmes. Aujourd'hui, les réseaux sociaux sont le plus grand
système de partage de mèmes dans notre société, ils ont souvent une vie
très courte mais une résonance globale, il suffit de penser à la chanson
``Gangnam Style''.


\subsubsection{Historique de l'objet}

Le meme a été théorisé par Richard Dawkins, biologiste, éthologiste et
théoricien de l'évolution, dans son essai ~``The Selfish Gene'' (Le Gène
égoïste). \cite{gene}

Le terme a pour origine les mots «~gène~» et «~mimesis~». En effet selon
R. Dawkins les cultures évoluent selon des lois analogues aux lois de
l'évolution de Darwin et Wallace, qui exerce une pression sur les gènes.
Les principes de variations et de sélection. \cite{larousse}

Le même et la théorie du mème (mémétique) s'inscrivent dans le
prolongement de la théorie de l'évolution a des objets non biologique,
dans le but de mieux comprendre le comportement des phénomenes

Il définit tout élément d'une culture ou d'un ensemble de comportements
répliqué et transmit d'un individu à un autre indépendamment de la
génétique. L'imitation par exemple, est le moyen privilégié. En d'autres
termes, c'est une idée qui a cours.

Etant défini comme l'unité de la culture, le mème n'a pas a proprement
parler d'origine. En effet le mème est aussi vieux que la culture
elle-même. Il a néanmoins connu de Nombreuses évolutions. A la fois en
ce qui concerne les supports de diffusions mais aussi dans leur portée.

A l'origine, les éléments culturels (les mèmes) était vraisemblablement,
communiqués par les gestes et la parole. Ils avaient donc une portée
locale, celle d'un groupe de personne réduit. L'arrivée de l'écriture
entraine la possibilité d'exporter les mèmes, d'échanger les cultures
plus loin et plus aisément, car les même peuvent ainsi durée plus
longtemps dans le temps qu'une simple parole. Enfin Internet et les
nouvelles technologies de communications rendent maximales la portée et
diversifie un peu plus les supports, tout en etant un moyen de partager
les mèmes de toutes sorte d manière efficace


\subsubsection{Aspect économique}

Il serait malaisé de définir un impact économique des mèmes. En effet il
ne s'agit pas d'un terme désignant un objet précis, mais d'un concept
très vaste comme on a pu l'évoquer plus tôt. On pourrait néanmoins se
pencher sur un cas particulier. Les «~mèmes~internet~» qui nous touchent
tous les jours sur les réseaux sociaux de l'internet, représentent une
manne financière importante quoi que difficile à exploiter.\cite{lorenz} En effet sur
internet les mèmes populaires ont une durée de vie très faible, ainsi il
serait difficile d'envisager une vente rentable de produit dérivée sur
une durée de seulement quelques semaines. S'ajoute à cela le caractère
imprévisible de la popularité des mèmes.\cite{reuter}


\subsubsection{Les diverses utilisations}

Le plus grand champ d'utilisation des mèmes et de la mémetique est dans
le secteur des études anthropologiques et sociologiques. En analyser le
fonctionnement permet aux experts de comprendre les ethnies et les
différentes cultures. Par exemple, Richard Dawkins, qui introduit le
terme, fait un grand excursus dans ``L'illusion de Dieu'' en expliquant
la présence et diffusion des idéaux religieux dans la société. Il
examine la mémétique et l'origine des mèmes ``divins'' dès l'antiquité
et leur parcours à travers l'humanité.

Le mème est amplement utilisé aujourd'hui dans le champ du
divertissement car tous les mouvements musicaux et artistiques dérivent
d'une transmission virale de certains mèmes. Tel comportement
reproductif explique l'instabilité générationnelle et la formation de
phénomènes très populaires mais ayant un arc temporel assez réduit. Cela
est hyperbolisé par l'avènement des réseaux sociaux qui sont un moyen
très efficace pour la transmission.

L'imitation, élément fondamental de la diffusion, est amplement utilisée
dans la politique, car les représentants veulent obtenir le consensus.
Cela est possible parmi un partage de valeurs, donc une exaspération des
mèmes qui sont le plus présents dans le pool mémétique local. Ce
phénomène peut avoir des résultats négatifs comme dans l'affaire italien
où le parti ``Lega'' a partagé les mèmes ``défectueux'' disant que les
vaccins provoquent l'autisme. L'utilisation médiatique de ce mème a eu
une résonance énorme dans la société, qui maintenant possède ce mème,
voir cette idée fausse. \cite{imitation}

Enfin, la compréhension des mèmes offre une vaste capacité publicitaire,
car plus une publicité utilise les mèmes actuellement très diffus, plus
elle a une chance d'avoir des effets sur les individus qui les possèdent
et qui ont une majeure compatibilité. Pour cette raison, l'objectif
médiatique est d'impliquer une grande partie des gens à travers une
résonance constructive entre le mèmes de la publicité et des individus.

\subsubsection{Les applications}


\myparagraph{\'{E}tudes antropologiques et sociologiques}
Les champs d'étude des mèmes se concentrent dans le secteur de la
sociologie qu'on définis ``étude scientifique des sociétés humaines et
des faits sociaux'' et ``étude des groupes humains qui exercent un
métier (sociologie rurale), qui professent une foi, manifestent des
croyances (sociologie religieuse), qui s'intéressent à un phénomène
culturel, artistique (par exemplesociologie de la littérature)''. Dans
ce cas la mémétique représente l'instrument adapte à comprendre le
passage des usages et des langages entre les différentes tribus et
cultures qui peuvent apparemment n'avoir pas aucune relation.

En considérant le métier de Richard Dawkins, introducteur du terme mème,
on en déduit l'application dans le secteur de l'anthropologie
évolutionniste. Il est définis comme ``l'étude de l'homme et des groupes
humains'' ou comme la ``théorie philosophique qui met l'homme au centre
de ses préoccupations''. \cite{gene} Puis on observe son spécialisation sur
l'évolutionnisme qui considère que toute culture est le résultat d'un
processus constant d'évolution. Cela nous souligne l'application du
concept des mèmes dans l'analyse du processus de passage et diffusion
des ``allèles'' mémétiques dans les cultures et dans l'histoire.


\myparagraph{Divertissement}
Chaque entité du champ sémantique du divertissement, dans ses aspect
sociaux d'amusement et de distraction, dépend directement des mèmes car
ceux-ci ont besoin d'être partagés pour survivre. Cette affirmation est
applicable dans la musique, surtout dans la composition de forme libre
ou de forme suite, dans le théâtre et chaque loisir, en particulier ceux
qui se développent comme activité d'équipe. Divertissement dérive du
latin \emph{divertere}, composé par \emph{di} signifiant ``de'' voir
éloignement et \emph{vertere} chargé du sens de ``tourner'', donc envers
un sentier autre part. \cite{IL} Dans l'ensemble la mot énonce une activité qui
procure plaisir plaisir et qui tient l'âme loin des préoccupations.


\myparagraph{Politique}
La politique est l'ensemble des options prises collectivement ou
individuellement par les gouvernants d'un État dans quelque domaine que
s'exerce leur autorité (domaine législatif, économique ou social,
relations extérieures) voire la méthode particulière de gouvernement,
manière de gouverner. La mémétique est appliqué dans la recherche du
\emph{consesus} à travers des sondages et des analyses statistiques sur
les intérêts et volontés des gens d'un secteur déterminé. Cette activité
permet de reconnaitre les mèmes membres d'une société, en suite ils
seront analysés et sélectionnes pour être assimilés par la figure du
parti ou du politicien. En effet, vu que il sont respectivement une
organisation structurée dont les membres mènent une action collective
dans la société aux fins de réaliser un programme politique et une
figure représentant une majorité du peuple, ils doivent partagés les
idéaux voir les mèmes avec la majeure partie.


\myparagraph{Publicité}
La publicité se manifeste comme une activité ayant pour but de faire
connaître une marque, d'inciter le public à acheter un produit, à
utiliser tel service, etc. ; ensemble des moyens et techniques employés
à cet effet (abréviation familière pub). Elle est concrètement un
annonce, un encart, un film, etc., conçus pour faire connaître et vanter
un produit, un service, etc. Dans ce domaine il s'applique le même
principe que pour la politique: il faut capturer le gents et cela est
possible seulement grâce à la correspondance dés mèmes car dans
l'histoire ancienne la reconnaissance tribale arrivait par la
comparaison des mèmes. \cite{dang:hal-00576540}  S'ils étaient suffisamment similaires, ça parait
évidant d'être du même groupe sociale, c'est le principe que permet
aujourd'hui à des italiens de se reconnaitre entre eux dans et hors de
leur pais natal. Enfin, la publicité veut produire le même effet pour
être les plus familial et efficace.


\subsection{Lieux}


\subsubsection{Lieux d'étude et conception}

L'étude des mème est surtout informative dans la mesure où elle nous
renseigne sur la manière dont se transmet l'information~: Afin de
déterminer quels supports sont les plus efficaces, quel sont les formes
les plus susceptibles d'être répliquées.

C'est dans ce but que Le NaN (Networks \& agents Network) un groupe du
CNetS (Center for complex Network and Systems Research, un département
de l'université de l'Indiana) étudie à travers divers projets, la
diffusion des mèmes, en particulier sur internet.

Le NaN est un groupe d'étude de 15 personnes (une seule équipe)~: un (1)
professeur d'informatique Filippo Menczer principal chercheur, deux (2)
chercheurs (post doctorat), et douze (12) élèves diplômés.

Le groupe a reçu plusieurs subventions d'importantes d'organisations
nationales américaines, à savoir la Fondation Nationale pour la Science,
et le DARPA, (« Agence pour les projets de recherche avancée de défense
»), d'une valeur de respectivement 1 million et 1,2 millions de dollars.
Ces subventions montrent que les applications de la mémetique
représentent un enjeu majeur pour les Etats-Unis d'Amérique.

Ces contributions permettent la mise en marche de projets d'envergure.

La première a permis la création de OSoMe (Observatory on Social Media),
une plateforme permettant de visualiser la diffusion des mèmes sur les
résaux sociaux en fonction du temps, à l'aide de cartes dynamique, de
courbes de tendances, ou de reconstitution de réseaux d'échanges de
mèmes.

La seconde quant à elle, finance une étude sur la propagation de
l'information sur internet via le projet COSINE (Cognitive Online
Simulation of Information Network Environments), qui vise à simuler un
cadre afin d'étudier les phénomènes sociaux a toutes échelles
géographiques et temporelles.


\subsubsection{Lieux de construction}
Le mème étant l'unité de la culture il n'exite pas a proprement parler
de lieu précis où sont construit les mèmes. Les mèmes se créent partout
où la culture peut se transmettre. On peut par contre évoquer des
support de construction.

Les œuvres artistiques constituent un support fréquemment à l'origine de
mèmes.

Artistes, phénomenes sociaux, théories scientifiques, traditions.


\subsubsection{Lieux d'utilisation}
Définir des lieux concrets d'application pour le mèmes est très
difficile car ils sont très diffus. Il suffit de considérer que tout
objet artificiel peut être considéré comme support mèmetique et chaque
humain est une véritable librairie de mèmes. Les mèmes sont partout et
en considérant l'utilisation primaire elle est conséquence de la
présence de tout mème, par contre une utilisation commerciale concentrée
dans tout centre urbain, surtout dans les métropoles.

\myparagraph{Pool et identité mèmetique}
Il est possible d'identifier différents niveaux de pools mèmétiques
parce que déjà un individu est considéré comme un pool. Par contre un
pool mèmétique est aussi un groupe de gens peu ou beaucoup entendu,
comme une bande, un centre urbain, une ethnie. \cite{imitation}

En observant deux individus ou deux groupes on peut en analyser les
similitudes mèmétiques en trouvant des mèmes analogues et d'autres
complètementdifférents. En faisant un parallélisme avec la génétique on
peut affirmer qu'une forte affinité entre les mèmes de deux individus
est conséquence d'uneappartenance au même groupe culturel.


\myparagraph{Hiérarchie des territoires et interactions}
Bien qu'on ne puisse pas établir un lieu physique aux mèmes, on en peut
analyser la concentration par rapport aux interactions. Aujourd'hui les
centres de positionnement et déplacement des mèmes correspond au lieux
plus développés et visités dans le monde. Des exemples sont des grandes
villes comme New York, Paris, London etc. ou bien des centres culturels
comme des musées ou lieux d'échanges.

Les territoires les plus avancées sont au centre des échanges et
interactions des mèmes parce que cela fonctionne par la conversation de
deux individus ou par la lecture d'une affiche ou encore d'une évènement
affolé. Ces phénomènes sont évidemment plus fréquentes dans les lieux ou
la population et les biens ce concentrent avec une rythme sûrement plus
rapide que dans les lieux isolées.


\myparagraph{Supports, technologies et déplacement humain}

\emph{\textbf{Globalisation}}: La globalisation a marqué un changement
radical dans le déplacement et diffusion des mèmes car ce phénomène a
abondement centralisé les mèmes en réduisant note vilement du quantité
une grand partie de ceuxqu'on peut définir ``minoritaires'' en
favorisant une énorme copie et distribution des mèmes globalisées.
L'exemple rampante est l'identité de la Coca Cola et les éléments
culturels en arrière plan.

\emph{\textbf{Réseaux sociaux}}: Les réseaux sociaux sont le plus grand
et plus vite bouleversement dans la transmission des mèmes. Ils ont
permis un échange instantané entre les gens sans limite de distance en
augmentant le rythme des échanges d'un facteur incalculable. \cite{gach:tel-01037937} De plus,
les réseaux sociaux ont provoqué la formation d'une série des mèmes qui
connaissent une distribution presque globale mais avec une vie très
courte, paradoxalement ce manifestent sous la forme des memes indendus
dans le sense comique.

\emph{\textbf{Livres}}: Les livres ce sont des véritables systèmes de
stockage de mèmes. Écrits dans une époque et un lieux bien définis, ils
deviennent un miroir des habitudes et façons de faire de ce lieux. Les
mèmes restent ``inactif'' dans les pages du livre, mais par exemple si
quelqu'un lit un vieux livre abandonné, y trouve une expression dépassée
et commence à la répéter, il fait ``revivre'' le mème en le disposant à
une possible nouvelle diffusion. \cite{life}

\emph{\textbf{Musique}}: La musique peut être rapprochée du livre parce
que elle est directement affecté d'une époque et d'un lieux bien
définis, mais la diffusion elle est souvent plus explosive et à part de
cas très limités aussi la durée tend à être limitée. De plus, la musique
emporte une grande quantité de mèmes sur les compétences, qualités
musicales et instruments, mèmes que peuvent être seulement transcrits du
point de vue théorique dans les livres.

Supports culturels

\subsubsection{Lieux d'application}

\emph{\textbf{Divertissement}}: C'est le grand domaine d'application des
mèmes, ils sont à la base de chaque mouvement ou action. Ce qui est
divertissant est définit par les mèmes représentantes le spectre d'une
culture. En fait cela explique pourquoi certaines choses amusent les uns
et pas les autres. Il faut considérer que souvent des individus
appartenant au même pool ont des mèmes liés au divertissements
similaires.

L'application de ce concept explique la diffusion géographique des
phénomènes de divertissement. Par exemple, les mèmes linguistiques sont
déterminants dans le positionnement des limites spatiales parce qu'elles
sont fortement caractéristiques des différents réalités locales. En
opposition, des êtres qui portent des mèmes proches aux instinct humains
sont plus favorable à une diffusion pas déterminée ou limitée dans les
territoires.

\emph{\textbf{Politique}}: Ici les enjeux des mèmes sont vitaux par
conséquent au concept de représentativité. Selon ce raisonnement les
élus devraient partager une vaste quantité des mèmes avec la plupart du
peuple. Cela veut signifier que souvent ils proviennent de la même
culture et bien du même lieux. Une analyse de ce phénomène correspond
aux données statistiques de provenance et de liaison à un lieux des
représentants.

Si on considère les campagnes politiques dans des réalités variées comme
celle italienne, on voit que parfois ils ont essayé d'envoyer des
façades différentes adaptées aux mèmes des petites réalités locales.
Enfin cela veut dire que afin d'analyser un facteur politique il faut
avoir de différents points de vue, déplacés par des lieux physiques et
culturels.

\emph{\textbf{Publicité}}: L'application dans le secteur de la publicité
est fortement dépendante des lieux parce que le taux d'efficacité est
conséquent à l'affinité des mèmes entre la publicité et le publique
concerné. Pour cette raison elle doit être pensée d'une façon spécifique
pour une certaine cible afin de afin de la toucher avec un efficaité
maximum et pas la grande masse avec une efficacité médiocre. Aujourd'hui
la technique dans le big data du profiling permet d'envoyer à une cible
bien précise et défini une pub qui le concerne et qui a une forte
similarité mémétique.

\subsection{Métiers}
\paragraph{Chercheur en science des résaux :} 
Consiste en l’élaboration de modèles pour  l’etude des interactions entre les individus ou objet, en l’occurrence (les flux d’information et de mèmes) il s’agit d’une approche plus sociologique.
Le métier touches principalement aux domaines de la science de l’information, et des sciences cognitives ou des résaux sociaux.
\paragraph{Sociologue :}
Consiste en la construction de modèles de fonctionnement de la société et des mécanismes sociaux. En l’occurrence les memes. Il peut donc s’agir d’étudier le langage ou les structures sociales autour de l’information.
Le métier touche aux disciplines de la philosophie et des sciences sociales
\subsection{Disciplines}

\begin{itemize}
\tightlist
\item
  Biologie
\item
  Sociologie
\item
  Philosophie
\item
	Informatique
\end{itemize}


\subsection{Parcours de formation}

Vu que il est possible d'étudier les mèmes tels quels seulement parmi leur expression phénotypique le parcours de formation proposé concerne activement le monde du numérique. Cette choix a été prise car les mèmes s'échangent d'une façon rapide, continue et surtout observable sur internet et sur les réseaux sociaux. 

Les études se concentrent sur l'informatique avec un intérêt pour la sociologie qui pourrait être suivie comme mineure ou comme parcours de formation individuel.
Puis, en se spécialisant dans le numérique il faudrait choisir de parcours et/ou masters liées à ces concepts: "Network  science,  Web  science,  data  science,  computational  social  science,  diffusion  of  information  and  
disinformation  in  social  media,  complex  networks  and  systems,  agent  based modeling,  applied  machine learning  and  data  mining,  cognitive  science".

Cela fournira une connaissance développée sur les réseaux afin d'en étudier les données. Mais ceci n'est pas suffisant: il serait bien d'ajouter des formations anthropologues/sociologues afin de comprendre la transmission et l'entité du mème pour comprendre et analyser les données obtenues.

\subsection{Conclusion}
Le mème est un objet fascinant dont on semble connaitre le fonctionnement, bien que celui demeure néanmoins difficile a reproduire « artificiellement » car imprévisible. Par ailleurs aux vu des enjeux naissant autour de cet objet, de plus de métiers vont venir éclaircir ces problématiques dans les années à venir. 

\newpage
\bibliographystyle{plain}
\bibliography{memoire}

\end{document}
